\section{Limitations and Threats to Validity}
\label{sec:limitations}

The technical evaluation establishes behavior for three purposively authored
development cases, not frequencies in a population of rooms, domains, or
authors.  Every natural asset has an explicit asset-level owner, so group and
zone fallback and ambiguity are exercised only in derived copies.  No
independently authored or held-out scene tests generalization.

The non-regression comparison is controlled rather than independent.  Both
representations originate from shared semantic artifacts and are classified by
one routing function after separate parsing; the direct-wiring projection also
supplies expected routes for the runtime sweep.  Agreement therefore detects
migration divergence but cannot validate the shared policy against an external
oracle.  Binding identifiers are checked for agreement within one pinned model,
not against independently authored labels.

Fault injection covers three common and three contract-only mutation types.
Detection and localization do not imply complete fault coverage or repair, and
scripted deltas are not developer effort.  The V2 assembler, validators, and
materializer share a codebase, so their agreement is not independent standards
conformance.  Runtime tests use a deterministic structured stub, isolating
admission, evidence, and rollback but not answer correctness or naturalness.
Timing covers only in-memory adapter work and excludes the end-to-end path.

Study U is an unmoderated, single-condition convenience study of one brand room.
The current cut contains 45 desktop/laptop and 20 tablet/smartphone sessions,
but no headset session.  Device heterogeneity, self-selection, and the young,
technically weighted sample limit generalization.  With 65 sessions, binary
proportions have at worst an approximately 12-percentage-point 95\% Wilson
half-width; this supports broad descriptive patterns, not fine subgroup or
population claims.  Without a baseline, observed comprehension cannot be
causally attributed to FunctionalMLDS rather than the interface and tasks.

The two user tasks were presented in a fixed order and both targeted successful
paths: one grounded object answer and one permitted direct handoff.  Study U
therefore does not directly test whether users understand the central
fail-closed states in Table~\ref{tab:failure-model}, including ambiguity,
unavailable capabilities, unreachable providers, or rollback.  Only 22
participants happened to rate an error explanation, so that result is
opportunistic and exploratory.  A controlled failure-comprehension task is
needed before claiming that the ``clarify or guide'' and rollback behavior is
intelligible to users.

The questionnaire records perceived handoffs but no linked per-session backend
handoff events, preventing individual log--report agreement.  The custom items
are not a validated psychometric scale, and the study does not evaluate
long-term use, trust calibration, workload, immersion, or developer diagnosis.
It tests only whether the current cues support target, answer, handoff, and
responsibility comprehension.

The project remains static within a session.  Drift is detected, but there is
no live migration for objects created, deleted, merged, or rebound during a
conversation.  Online routing permits at most one direct handoff per request;
transitive reachability is an offline diagnostic.  Concurrent writes and
simultaneous visitors were not tested.  Finally, the selected EAST-ADL subset is
a traceability foundation; full conformance has not been independently assessed.

\section{Conclusion}
\label{sec:conclusion}

We presented an executable interaction contract that keeps a selected spatial
entity, responsible embodied agent, declared capability, concrete backend
action, and returned evidence within one pinned model version.  The client
makes unresolved selection explicit, while the backend re-resolves identity,
responsibility, route, and action before generation and rolls back
response-dependent violations.

Across three authored environments, all 93 assets and 186 target-specific
declaration pairs were complete.  Migration preserved all 93 owner decisions
and 459 route classifications under a shared policy.  The runtime admitted all
298 local or direct routes with matching evidence and rejected all 161
transitive or unreachable routes before the response stub without retained
mutation.  A strict provider contract additionally detected and localized all
nine injected cross-layer faults that rely on relations absent from the direct
artifacts.

The complementary user study shows why technical enforcement and interface
intelligibility must be evaluated separately.  Participants generally reported
clear target--answer fit and visible handoffs, but only 37/65 identified the
intended object-and-topic responsibility rule, and the usefulness of several
specialized agents received mixed ratings.  These results provide initial
evidence for successful-path intelligibility, not for comprehension of blocked
or rolled-back interactions.  FunctionalMLDS provides the evaluated fail-closed
control and evidence boundary, while the interface still needs clearer
responsibility and failure communication.
